\section{System Design}

\subsection{Simulator-owned legality}

The environment boundary is intentionally asymmetric. The simulator owns exact
rules, random outcomes, legal menus, transitions, and terminal conditions. The
controller owns representation and ranking. At every factorized stage the
wrapper requests the current legal list, encodes every option, selects one
current index, and submits that index. If model load or inference fails, a
seeded fallback still chooses from the supplied list. This cannot guarantee a
strategically good action, but it preserves the basic legality invariant
\begin{equation}
  a_{t,k}\in\mathcal{A}(I_{t,k}) \quad \text{for every executed stage }k.
\end{equation}

\subsection{Causal state and option encoding}

The base model uses a Transformer-style causal representation
\cite{vaswani2017attention}. Spatial board information, acting-seat private
state, public flags and counts, and realized action history are encoded without
future tokens. Each legal option receives structured identity, category,
source, destination, target, count, effect, and continuation context. A
factorized decoder compares the current legal set in parallel, then emits one
option at a time. For multi-select prompts, already selected options and the
remaining budget are part of the next stage; an explicit \textsc{stop} becomes
selectable only when the engine exposes it.

The score for current option $i$ can be written schematically as
\begin{equation}
  \ell_i = \ell_i^{\mathrm{base}}
    + g_{\mathrm{rule}}\,r_i
    + \sum_{j\in\mathcal{H}} m_{j,i}\,\rho_{j,i}\,\Delta_{j,i},
  \label{eq:fused-score}
\end{equation}
where $r_i$ is the rule-semantic residual, $g_{\mathrm{rule}}$ is its gated
activation, and each auxiliary route has an availability mask $m_{j,i}$,
bounded reliability $\rho_{j,i}$, and option-conditioned contribution
$\Delta_{j,i}$. Equation~\eqref{eq:fused-score} describes the architectural
contract, not a claim that every physical route is active. The census accepted
only the public-rule projection; unsupported branches keep exact-zero gates.

\subsection{Public-rule semantic repair}

The legacy card representation is preserved bit-for-bit. The repair is a
parallel residual whose schema includes:

\begin{itemize}
\item exact NUMBER values rather than a coarse bucket or text hash;
\item selection context, effect, context card, and option count;
\item minimum, maximum, selected, and remaining selection counts;
\item remaining damage and energy budgets when publicly defined;
\item stable source, target, area, slot, tool, and energy bindings;
\item public current/maximum HP, evolution stack, appearance-this-turn, bench
maximum, typed energy units, and once-per-turn flags when retained and proven;
\item pinned public catalog metadata through a separately gated adapter.
\end{itemize}

Global serials and incidental menu order are not semantic identities. Stable
within-observation entity bindings are preferred; ordinal position is a
last-resort fallback only when the simulator distinguishes otherwise identical
payloads. Text hashes may remain in the legacy card path but never serve as
exact mechanics authority for the new residual. The new path is initialized to
zero, and a layer-off test requires bit-identical base logits. This gives the
migration an executable rollback rather than a prose promise.

\subsection{Corpus-driven branch adjudication}

Implementation is not evidence of usefulness. Every proposed branch is first
physical, default-zero, and inert. The frozen schema is then applied to the
complete eligible corpus, and a census measures public semantic collisions,
selected-action involvement, deck and matchup distribution, transition
witnesses, and harmless equivalences. Only a branch supported by the census may
enter training or action scoring.

For the current snapshot, adjudication accepted
\texttt{public\_rule\_semantic\_projection}. The public metadata residual,
repaired auxiliary heads, and all eight checklist-provenance gates remained
unsupported or unproven and therefore exact-zero. This distinction prevents a
paper from describing a head inventory as realized decision influence.

\subsection{Checklist provenance}

The schema records eight public questions at every actor-visible decision or
forced selection: current knockout hand threshold; cards safely spendable above
that threshold; replacement Alakazam line; unavoidable draws before the next
attack; bench Prize exposure; immediate disruption outcome; a robust line under
unknown Prize identity; and whether the game ends before a forced draw. Each
channel records factual evidence, availability, an unavailable reason,
per-option signed evidence when defined, the gate, and final residual. Unknown
or hidden evidence yields zero. The current corpus did not authorize the
checklist gate branch, so these records support trace and future calibration,
not a present action-influence claim.

\subsection{Typed auxiliary separation}

The r274-compatible policy structure contains existing policy, value,
strategic-head, Fusion, own-deck, and matchup paths. Revision 10 mapped immutable
r195-compatible weights into that structure and optimized every
architecture-present non-combo model weight. Combo tensors were retained only
for checkpoint and pack ABI compatibility, with zero loss, no optimizer
membership, and a disabled runtime route. The preserved Card2Vec branch is a
specific non-regression boundary: the semantic residual may accompany it but
cannot replace, re-embed, or rewire it.

Matchup routing is also causal. A route must resolve from supported public
identity; unknown, ambiguous, retired, or unsupported identities take the
base-policy bypass. Model materialization alone is not route selection, and
route selection alone is not evidence of improved outcomes.
